\chapter{Phases Of Gauge Theories}
\label{ch:global-phases-of-gauge-theories}

\ConventionsMixed

This chapter defines phases of gauge theories through infinite-volume QFT
data: states, correlators, extended operators, symmetry realization, and
limits of renormalized line observables.  Local gauge-variant fields do not
label phases; gauge-invariant operator data do.

\section{Phase Data}

Fix a class of regulators and a set of renormalized observables
\[
  \mathfrak O=\{\text{local operators, line operators, surface operators,
  background-field responses}\}.
\]
A phase is a connected region in the parameter space of infinite-volume
limits for which the following data vary continuously in the chosen topology:
the vacuum sector, the mass-gap or gapless spectral behavior, the correlation
lengths of local gauge-invariant operators, the expectation values and fusion
rules of extended operators, and the realization of ordinary and higher-form
symmetries.  A phase boundary is a locus where at least one of these data
ceases to extend continuously.

\section{Static Potentials From Wilson Loops}

Let \(W_R(C)\) be a renormalized Wilson loop in representation \(R\) of the
gauge group.  For a rectangle \(C_{T,L}\) of temporal length \(T\) and spatial
length \(L\), define the static potential by
\[
  V_R(L)=
  -\lim_{T\to\infty}\frac{1}{T}\log
  \left\langle W_R(C_{T,L})\right\rangle ,
\]
when the limit exists after perimeter renormalization.  A confining response
for the external charge \(R\) is
\[
  V_R(L)=\sigma_R L+o(L),\qquad \sigma_R>0 .
\]
A screened response is bounded as \(L\to\infty\).  A Coulomb response in four
spacetime dimensions has leading behavior
\[
  V_R(L)=V_\infty-\frac{\kappa_R}{L}+o(L^{-1})
\]
in the sector where massless gauge fields remain in the infrared.

\section{Higher-Form Symmetry Realization}

For a pure gauge theory with center one-form symmetry \(A\), a Wilson line
charged under \(A\) is an order operator for that symmetry.  In a gapped
phase, an area law for all nontrivially charged Wilson loops is the
confining realization of the electric one-form symmetry.  A perimeter law
for a charged Wilson loop is the broken realization.  Magnetic one-form
symmetries and 't Hooft loops give the dual diagnostics, with their charge
lattice determined by the global form of the gauge group and by discrete
theta terms.

\section{Higgs, Coulomb, And Confining Regions}

Suppose a gauge theory contains matter fields whose gauge-invariant composites
have expectation values that change the spectrum of charged excitations.  A
Higgs region is described by a mass gap for gauge bosons in a weakly coupled
coordinate chart and by screened Wilson loops for charges carried by the
matter sector.  A Coulomb region contains massless gauge fields and long-range
electric or magnetic fields.  A confining region has a positive string
tension for external charges that cannot be screened by dynamical matter.
These are descriptions of responses in a specified observable set; the phase
classification is the classification of the full infinite-volume data in
Section~1.

\section{Analytic Connectivity In Gauge-Higgs Models}

Consider a lattice gauge theory with compact gauge group and matter in a
representation that screens the fundamental Wilson charge.  At finite lattice
volume the partition function is an analytic function of the couplings.  A
cluster expansion gives open regions in the strong-gauge and large-Higgs
coupling domains where connected correlations of local gauge-invariant
operators decay exponentially and are analytic in the couplings.

\begin{proposition}[Gauge-Higgs analytic corridor]
\label{prop:gauge-higgs-analytic-corridor}
Assume the cluster expansions in the strong-gauge and large-Higgs domains
converge in a common connected corridor of lattice couplings, uniformly in
volume for local gauge-invariant observables.  Then those local observables
have a single analytic infinite-volume continuation throughout the corridor.
\end{proposition}

\begin{proof}
Uniform convergence of the cluster expansion gives an absolutely convergent
series for each connected local correlator with coefficients analytic in the
couplings.  The thermodynamic limit may therefore be taken term by term in
the corridor.  Analytic continuation along any path inside the connected
corridor gives the same value because the local correlator is represented by
one convergent analytic series on overlapping neighborhoods.  Hence the two
coordinate descriptions belong to one local-observable phase in the declared
observable set.
\end{proof}

The proposition concerns local gauge-invariant observables in the specified
corridor.  Extended probes, boundary conditions, higher-form backgrounds, and
heavy external charges may carry additional phase data and must be included
explicitly when they are part of \(\mathfrak O\).

\section{Topological Order And Deconfined Gauge Theory}

A gapped gauge theory can have a finite-dimensional Hilbert space on spatial
manifolds with nontrivial topology, nontrivial line-operator braiding, and
surface operators whose correlation functions depend on homology classes.
These data form the long-distance topological sector.  In continuum language
they are encoded by an infrared TQFT together with the map from microscopic
extended operators to topological operators.  This map is part of the phase
definition, because different microscopic theories may flow to the same
topological sector while differing in massive local spectra and symmetry
responses.

\begin{openproblem}[Complete phase invariant for gauge theories]
\label{op:complete-phase-invariant-gauge-theories}
Define a tractable invariant of gauge-theory phases that includes local
correlators, extended-operator categories, higher-form symmetry realization,
background responses, and massive spectra, and prove its stability under
finite local deformations inside a gapped phase.
\end{openproblem}
